\chapter{Schwache und elektroschwache Wechselwirkung}

\section{Ziel dieses Kapitels}

Dieses Kapitel behandelt die schwache Wechselwirkung und ihre Verbindung mit der
elektromagnetischen Wechselwirkung zur elektroschwachen Wechselwirkung.

Die schwache Wechselwirkung ist besonders wichtig fuer Beta-Zerfaelle, Neutrinos,
Flavorwechsel, Paritaetsverletzung und die Rolle der $W$- und $Z$-Bosonen.

\begin{notebox}
Dieses Kapitel behandelt die Stichwoerter:
\begin{itemize}
    \item schwache Wechselwirkung,
    \item geladener schwacher Strom,
    \item neutraler schwacher Strom,
    \item $W^+$-, $W^-$- und $Z^0$-Bosonen,
    \item Beta-Zerfall,
    \item Neutrinos,
    \item Flavorwechsel,
    \item Paritaetsverletzung,
    \item Fermi-Theorie,
    \item kurze Reichweite der schwachen Wechselwirkung,
    \item elektroschwache Vereinheitlichung,
    \item Photon und $Z^0$,
    \item Higgs-Mechanismus,
    \item Masse der schwachen Eichbosonen.
\end{itemize}
\end{notebox}

\section{Grundidee der schwachen Wechselwirkung}

Die schwache Wechselwirkung ist eine der fundamentalen Wechselwirkungen.
Sie wirkt auf Quarks und Leptonen.

\begin{definitionbox}
Die \emph{schwache Wechselwirkung} ist die fundamentale Wechselwirkung, die unter
anderem Beta-Zerfaelle, Neutrinoreaktionen und Quark-Flavorwechsel ermoeglicht.
\end{definitionbox}

\begin{conceptbox}
Typische Prozesse der schwachen Wechselwirkung sind:
\begin{itemize}
    \item Beta-Minus-Zerfall,
    \item Beta-Plus-Zerfall,
    \item Elektroneneinfang,
    \item Myonzerfall,
    \item Pionzerfall,
    \item Neutrinostreuung,
    \item Zerfaelle seltsamer Teilchen.
\end{itemize}
\end{conceptbox}

\begin{examtipbox}
Wenn in einem Prozess ein Neutrino vorkommt, ist fast immer die schwache
Wechselwirkung beteiligt.

Wenn sich ein Quark-Flavor aendert, ist ebenfalls die schwache Wechselwirkung
beteiligt.
\end{examtipbox}

\section{Reichweite und Staerke}

Die schwache Wechselwirkung hat eine sehr kurze Reichweite.
Das liegt daran, dass ihre Austauschteilchen massiv sind.

\begin{notebox}
Die Grundkraefte-Tabelle der Vorlesung ordnet die schwache Wechselwirkung als sehr
kurzreichweitig ein:
\[
    R_{\mathrm{schwach}} < 10^{-15}\,\mathrm{m}.
\]

Ihre relative Staerke ist deutlich kleiner als die der starken Wechselwirkung.
\end{notebox}

\begin{conceptbox}
Die Austauschteilchen der schwachen Wechselwirkung sind:
\[
    W^+,
    \qquad
    W^-,
    \qquad
    Z^0.
\]

Diese Bosonen sind sehr massiv.
Deshalb ist die Reichweite der schwachen Wechselwirkung sehr kurz.
\end{conceptbox}

\begin{examtipbox}
Kurze Reichweite der schwachen Wechselwirkung:
\[
    \text{massive Austauschteilchen}
    \quad\Rightarrow\quad
    \text{kurze Reichweite}.
\]
\end{examtipbox}

\section{Austauschteilchen der schwachen Wechselwirkung}

\begin{definitionbox}
Die schwache Wechselwirkung wird durch die schweren Eichbosonen
\[
    W^+,\quad W^-,\quad Z^0
\]
vermittelt.
\end{definitionbox}

\begin{conceptbox}
Es gibt zwei Arten schwacher Wechselwirkung:
\begin{itemize}
    \item geladener schwacher Strom mit $W^+$ und $W^-$,
    \item neutraler schwacher Strom mit $Z^0$.
\end{itemize}
\end{conceptbox}

\begin{center}
\begin{tabular}{|p{4.0cm}|p{4.0cm}|p{5.2cm}|}
\hline
Austauschteilchen & Ladung & typische Rolle \\
\hline
$W^+$ & $+1$ & geladener schwacher Strom \\
\hline
$W^-$ & $-1$ & geladener schwacher Strom \\
\hline
$Z^0$ & $0$ & neutraler schwacher Strom \\
\hline
\end{tabular}
\end{center}

\begin{examtipbox}
$W^\pm$ koennen die elektrische Ladung eines Fermions aendern.

$Z^0$ ist neutral und aendert die elektrische Ladung nicht.
\end{examtipbox}

\section{Geladener schwacher Strom}

Der geladene schwache Strom wird durch $W^+$- und $W^-$-Bosonen vermittelt.

\begin{definitionbox}
Ein \emph{geladener schwacher Strom} ist eine schwache Wechselwirkung, bei der ein
$W^+$- oder $W^-$-Boson ausgetauscht wird.
\end{definitionbox}

\begin{conceptbox}
Beim geladenen schwachen Strom kann ein Fermion in seinen Partner innerhalb eines
schwachen Dubletts uebergehen.

Beispiele:
\[
    d \rightarrow u + W^-,
\]
\[
    u \rightarrow d + W^+,
\]
\[
    e^- \rightarrow \nu_e + W^-.
\]
\end{conceptbox}

\begin{examplebox}
Beim Vertex
\[
    d \rightarrow u + W^-
\]
ist die elektrische Ladung erhalten:
\[
    -\frac{1}{3}
    =
    +\frac{2}{3}
    -
    1.
\]
\end{examplebox}

\begin{examtipbox}
Geladener schwacher Strom:
\begin{itemize}
    \item $W^\pm$ beteiligt,
    \item elektrische Ladung des Fermions kann sich aendern,
    \item Quark-Flavor kann sich aendern,
    \item Neutrinos koennen beteiligt sein.
\end{itemize}
\end{examtipbox}

\section{Neutraler schwacher Strom}

Der neutrale schwache Strom wird durch das $Z^0$-Boson vermittelt.

\begin{definitionbox}
Ein \emph{neutraler schwacher Strom} ist eine schwache Wechselwirkung, bei der ein
$Z^0$-Boson ausgetauscht wird.
\end{definitionbox}

\begin{conceptbox}
Das $Z^0$-Boson ist elektrisch neutral.
Deshalb bleibt die elektrische Ladung des Fermions am Vertex gleich.

Typische Struktur:
\[
    f \rightarrow f + Z^0.
\]
\end{conceptbox}

\begin{examplebox}
Ein neutrino-elektronischer Streuprozess kann ueber ein $Z^0$ laufen:
\[
    \nu_\mu + e^- \rightarrow \nu_\mu + e^-.
\]

Dabei bleiben die Teilchenarten am neutralen schwachen Vertex erhalten.
\end{examplebox}

\begin{examtipbox}
Neutraler schwacher Strom:
\begin{itemize}
    \item $Z^0$ beteiligt,
    \item keine Ladungsaenderung,
    \item kein geladener Leptonpartner muss entstehen,
    \item wichtig fuer Neutrinostreuung.
\end{itemize}
\end{examtipbox}

\section{Schwache Dubletts}

In der schwachen Wechselwirkung treten linkshaendige Fermionen in Dubletts auf.

\begin{conceptbox}
Lepton-Dubletts:
\[
    \begin{pmatrix}
        \nu_e \\
        e^-
    \end{pmatrix},
    \qquad
    \begin{pmatrix}
        \nu_\mu \\
        \mu^-
    \end{pmatrix},
    \qquad
    \begin{pmatrix}
        \nu_\tau \\
        \tau^-
    \end{pmatrix}.
\]
\end{conceptbox}

\begin{conceptbox}
Quark-Dubletts in einfacher Darstellung:
\[
    \begin{pmatrix}
        u \\
        d
    \end{pmatrix},
    \qquad
    \begin{pmatrix}
        c \\
        s
    \end{pmatrix},
    \qquad
    \begin{pmatrix}
        t \\
        b
    \end{pmatrix}.
\]
\end{conceptbox}

\begin{notebox}
In der vollstaendigen Theorie sind die schwachen Quarkzustande Mischungen der
Massen-Eigenzustaende.
Diese Mischung wird durch die CKM-Matrix beschrieben.
Fuer die Grundidee reicht:
$W^\pm$ koppeln up-artige Quarks an down-artige Quarks.
\end{notebox}

\section{Beta-Minus-Zerfall}

Der Beta-Minus-Zerfall ist ein klassisches Beispiel der schwachen Wechselwirkung.

\begin{formulabox}
Auf Nukleonenebene:
\[
    n \rightarrow p + e^- + \overline{\nu}_e.
\]
\end{formulabox}

\begin{formulabox}
Auf Quarkebene:
\[
    d \rightarrow u + W^-,
\]
\[
    W^- \rightarrow e^- + \overline{\nu}_e.
\]
\end{formulabox}

\begin{conceptbox}
Das Neutron hat Quarkinhalt:
\[
    n=udd.
\]

Das Proton hat Quarkinhalt:
\[
    p=uud.
\]

Beim Beta-Minus-Zerfall wird eines der Down-Quarks in ein Up-Quark umgewandelt.
\end{conceptbox}

\begin{examtipbox}
Beta-Minus-Zerfall:
\[
    n \rightarrow p + e^- + \overline{\nu}_e.
\]

Merken:
\[
    d \rightarrow u + W^-.
\]
\end{examtipbox}

\section{Beta-Plus-Zerfall}

Beim Beta-Plus-Zerfall wird ein Proton in ein Neutron umgewandelt.

\begin{formulabox}
Auf Nukleonenebene:
\[
    p \rightarrow n + e^+ + \nu_e.
\]
\end{formulabox}

\begin{formulabox}
Auf Quarkebene:
\[
    u \rightarrow d + W^+,
\]
\[
    W^+ \rightarrow e^+ + \nu_e.
\]
\end{formulabox}

\begin{conceptbox}
Das Proton hat Quarkinhalt:
\[
    p=uud.
\]

Das Neutron hat Quarkinhalt:
\[
    n=udd.
\]

Beim Beta-Plus-Zerfall wird eines der Up-Quarks in ein Down-Quark umgewandelt.
\end{conceptbox}

\begin{examtipbox}
Beta-Plus-Zerfall:
\[
    p \rightarrow n + e^+ + \nu_e.
\]

Merken:
\[
    u \rightarrow d + W^+.
\]
\end{examtipbox}

\section{Elektroneneinfang}

Beim Elektroneneinfang wird ein Elektron aus der Atomhuelle vom Kern eingefangen.

\begin{formulabox}
Auf Nukleonenebene:
\[
    p + e^- \rightarrow n + \nu_e.
\]
\end{formulabox}

\begin{formulabox}
Auf Quarkebene:
\[
    u + e^- \rightarrow d + \nu_e.
\]
\end{formulabox}

\begin{conceptbox}
Im Kern wird ein Proton in ein Neutron umgewandelt.
Die Massenzahl bleibt gleich, aber die Kernladungszahl nimmt um eins ab:
\[
    A=\text{konstant},
    \qquad
    Z \rightarrow Z-1.
\]
\end{conceptbox}

\begin{examtipbox}
Elektroneneinfang:
\[
    p+e^- \rightarrow n+\nu_e.
\]

Im Nuklid:
\[
    Z \rightarrow Z-1,
    \qquad
    A=\text{konstant}.
\]
\end{examtipbox}

\section{Myonzerfall}

Der Myonzerfall ist ein rein leptonischer schwacher Zerfall.

\begin{formulabox}
\[
    \mu^- \rightarrow e^- + \overline{\nu}_e + \nu_\mu.
\]
\end{formulabox}

\begin{conceptbox}
Man kann ihn ueber ein virtuelles $W^-$ verstehen:
\[
    \mu^- \rightarrow \nu_\mu + W^-,
\]
\[
    W^- \rightarrow e^- + \overline{\nu}_e.
\]
\end{conceptbox}

\begin{examplebox}
Leptonenzahl:
\[
    L_\mu:
    \qquad
    1 \rightarrow 1.
\]

Elektron-Leptonenzahl:
\[
    L_e:
    \qquad
    0 \rightarrow 1-1=0.
\]

Die Leptonenzahlen sind in dieser einfachen Beschreibung erhalten.
\end{examplebox}

\section{Neutrinos}

Neutrinos sind elektrisch neutrale Leptonen.
Sie nehmen an der schwachen Wechselwirkung teil.

\begin{definitionbox}
\emph{Neutrinos} sind elektrisch neutrale Leptonen, die ueber die schwache
Wechselwirkung wechselwirken.
\end{definitionbox}

\begin{conceptbox}
Es gibt drei Neutrino-Flavors:
\[
    \nu_e,
    \qquad
    \nu_\mu,
    \qquad
    \nu_\tau.
\]

Dazu gehoeren die Antineutrinos:
\[
    \overline{\nu}_e,
    \qquad
    \overline{\nu}_\mu,
    \qquad
    \overline{\nu}_\tau.
\]
\end{conceptbox}

\begin{notebox}
Neutrinos sind schwer nachzuweisen, weil sie keine elektrische Ladung tragen und
nicht stark wechselwirken.

Man weist sie ueber schwache Reaktionen nach.
\end{notebox}

\begin{examtipbox}
Neutrino im Prozess:
\[
    \Rightarrow
    \text{schwache Wechselwirkung}.
\]
\end{examtipbox}

\section{Neutrinonachweis durch inverse Beta-Reaktion}

Eine wichtige Nachweisreaktion ist die inverse Beta-Reaktion.

\begin{formulabox}
\[
    \overline{\nu}_e + p \rightarrow n + e^+.
\]
\end{formulabox}

\begin{conceptbox}
Ein Elektron-Antineutrino trifft auf ein Proton.
Dabei entstehen ein Neutron und ein Positron.

Das Positron kann anschliessend mit einem Elektron annihilieren:
\[
    e^+ + e^- \rightarrow \gamma+\gamma.
\]
Diese Photonen koennen im Detektor gemessen werden.
\end{conceptbox}

\begin{examtipbox}
Inverse Beta-Reaktion:
\[
    \overline{\nu}_e + p \rightarrow n + e^+.
\]

Sie ist ein Standardbeispiel fuer Neutrinodetektion.
\end{examtipbox}

\section{Flavorwechsel in der schwachen Wechselwirkung}

Die schwache Wechselwirkung kann Quark-Flavors aendern.

\begin{conceptbox}
Starke und elektromagnetische Wechselwirkungen erhalten Quark-Flavor.

Die schwache Wechselwirkung kann Quark-Flavor aendern:
\[
    d \rightarrow u,
\]
\[
    s \rightarrow u,
\]
\[
    b \rightarrow c.
\]
\end{conceptbox}

\begin{notebox}
Deshalb koennen seltsame Teilchen stark erzeugt, aber schwach zerfallen.

Starke Erzeugung:
Strangeness wird paarweise erzeugt und insgesamt erhalten.

Schwacher Zerfall:
Strangeness kann sich aendern.
\end{notebox}

\begin{examtipbox}
Wenn Strangeness in einem Zerfall nicht erhalten ist, ist das ein Hinweis auf
schwache Wechselwirkung.
\end{examtipbox}

\section{CKM-Mischung}

Bei Quarks koppelt das $W$-Boson nicht nur streng innerhalb einer Generation.
Es gibt Mischungen zwischen den Generationen.

\begin{definitionbox}
Die \emph{CKM-Matrix} beschreibt die Mischung der Quark-Flavors in geladenen
schwachen Wechselwirkungen.
\end{definitionbox}

\begin{conceptbox}
Ein up-artiges Quark kann ueber ein $W$-Boson an verschiedene down-artige Quarks
koppeln.

Zum Beispiel kann ein Charm-Quark schwach zerfallen in:
\[
    c \rightarrow s + W^+,
\]
oder unterdrueckt auch in:
\[
    c \rightarrow d + W^+.
\]
\end{conceptbox}

\begin{notebox}
Fuer viele Aufgaben in dieser Vorlesung reicht die qualitative Aussage:
Die CKM-Matrix erlaubt und gewichtet Flavorwechsel zwischen Quarkgenerationen.
\end{notebox}

\begin{examtipbox}
CKM-Matrix:
\[
    \text{Quarkmischung in der geladenen schwachen Wechselwirkung}.
\]
\end{examtipbox}

\section{Paritaetsverletzung}

Eine der wichtigsten Eigenschaften der schwachen Wechselwirkung ist die Verletzung
der Paritaet.

\begin{definitionbox}
\emph{Paritaet} beschreibt das Verhalten eines Systems unter Raumspiegelung:
\[
    \vec r \rightarrow -\vec r.
\]
\end{definitionbox}

\begin{conceptbox}
Starke und elektromagnetische Wechselwirkungen erhalten Paritaet.

Die schwache Wechselwirkung verletzt Paritaet.
Das bedeutet:
Ein schwacher Prozess und sein gespiegelter Prozess sind nicht gleich wahrscheinlich.
\end{conceptbox}

\begin{examtipbox}
Paritaetsverletzung ist ein Kennzeichen der schwachen Wechselwirkung.
\end{examtipbox}

\section{Chiralitaet und Linkshaendigkeit}

Die schwache Wechselwirkung koppelt im Standardmodell besonders an linkshaendige
Fermionen und rechtshaendige Antifermionen.

\begin{definitionbox}
\emph{Chiralitaet} ist eine relativistische Eigenschaft von Fermionfeldern.
Sie ist eng mit der Frage verbunden, ob ein Teilchen links- oder rechtshaendig ist.
\end{definitionbox}

\begin{conceptbox}
Geladene schwache Wechselwirkung koppelt an:
\begin{itemize}
    \item linkshaendige Fermionen,
    \item rechtshaendige Antifermionen.
\end{itemize}

Diese Asymmetrie ist der Grund fuer maximale Paritaetsverletzung in geladenen
schwachen Prozessen.
\end{conceptbox}

\begin{notebox}
Fuer masselose Teilchen fallen Chiralitaet und Helizitaet zusammen.
Bei massiven Teilchen muss man genauer unterscheiden.
Fuer die Grundidee reicht:
Die schwache Wechselwirkung behandelt links und rechts nicht gleich.
\end{notebox}

\section{Fermi-Theorie}

Vor der Entdeckung der $W$- und $Z$-Bosonen wurde der Beta-Zerfall durch eine
Kontaktwechselwirkung beschrieben.

\begin{definitionbox}
Die \emph{Fermi-Theorie} beschreibt schwache Prozesse bei niedrigen Energien als
Vier-Fermion-Kontaktwechselwirkung.
\end{definitionbox}

\begin{conceptbox}
Im Beta-Minus-Zerfall koppeln effektiv vier Fermionen an einem Punkt:
\[
    n \rightarrow p+e^-+\overline{\nu}_e.
\]

Auf Quarkebene:
\[
    d \rightarrow u+e^-+\overline{\nu}_e.
\]
\end{conceptbox}

\begin{formulabox}
Die Staerke der niederenergetischen schwachen Wechselwirkung wird durch die
Fermi-Konstante beschrieben:
\[
    G_F.
\]
\end{formulabox}

\begin{notebox}
Die moderne Theorie erklaert die Kontaktwechselwirkung als Naeherung fuer den
Austausch eines sehr schweren $W$-Bosons bei niedrigen Energien.
\end{notebox}

\begin{examtipbox}
Fermi-Theorie:
\[
    \text{niedrige Energie}
    \quad\Rightarrow\quad
    \text{Kontaktwechselwirkung}.
\]

Moderne Sicht:
\[
    \text{Kontaktwechselwirkung}
    \quad\approx\quad
    \text{Austausch eines schweren } W\text{-Bosons}.
\]
\end{examtipbox}

\section{Warum ist die schwache Wechselwirkung schwach?}

Bei niedrigen Energien erscheint die schwache Wechselwirkung schwach, weil die
Austauschteilchen sehr schwer sind.

\begin{conceptbox}
Ein Prozess mit virtuellem $W$-Boson enthaelt einen Propagator.
Bei niedrigen Impulsuebertraegen ist der Effekt des schweren $W$-Bosons stark
unterdrueckt.

Anschaulich:
Ein sehr schweres Austauschteilchen erzeugt eine sehr kurze Reichweite und eine
kleine effektive Kopplung bei niedrigen Energien.
\end{conceptbox}

\begin{notebox}
Bei Energien in der Groessenordnung der $W$- und $Z$-Massen wird sichtbar, dass
schwache und elektromagnetische Wechselwirkung Teile einer gemeinsamen
elektroschwachen Theorie sind.
\end{notebox}

\section{Elektroschwache Vereinheitlichung}

Die elektromagnetische und die schwache Wechselwirkung sind bei hohen Energien
nicht unabhaengig, sondern werden durch eine gemeinsame Theorie beschrieben.

\begin{definitionbox}
Die \emph{elektroschwache Wechselwirkung} ist die vereinheitlichte Beschreibung der
elektromagnetischen und schwachen Wechselwirkung.
\end{definitionbox}

\begin{conceptbox}
In der elektroschwachen Theorie entstehen die beobachteten Bosonen
\[
    \gamma,
    \qquad
    W^+,
    \qquad
    W^-,
    \qquad
    Z^0
\]
aus gemeinsamen Eichfeldern.
\end{conceptbox}

\begin{examtipbox}
Elektroschwache Vereinheitlichung bedeutet:
Photon, $W^\pm$ und $Z^0$ gehoeren zu einer gemeinsamen theoretischen Struktur.

Bei niedrigen Energien erscheinen elektromagnetische und schwache Wechselwirkung
sehr verschieden.
\end{examtipbox}

\section{Photon und Z-Boson}

Photon und $Z^0$ sind beide neutrale Eichbosonen, aber sie haben unterschiedliche
Eigenschaften.

\begin{center}
\begin{tabular}{|p{4.0cm}|p{4.5cm}|p{4.5cm}|}
\hline
Eigenschaft & Photon $\gamma$ & $Z^0$-Boson \\
\hline
Wechselwirkung & elektromagnetisch & schwach neutral \\
\hline
elektrische Ladung & $0$ & $0$ \\
\hline
Masse & $0$ & gross \\
\hline
Reichweite & unendlich & sehr kurz \\
\hline
koppelt an & elektrische Ladung & schwache Ladungen \\
\hline
\end{tabular}
\end{center}

\begin{conceptbox}
Das Photon ist masselos.
Deshalb hat die elektromagnetische Wechselwirkung eine unendliche Reichweite.

Das $Z^0$ ist massiv.
Deshalb ist die neutrale schwache Wechselwirkung kurzreichweitig.
\end{conceptbox}

\section{W-Bosonen}

Die $W$-Bosonen sind elektrisch geladen.

\begin{definitionbox}
Die $W$-Bosonen sind die geladenen Austauschteilchen der schwachen Wechselwirkung:
\[
    W^+,
    \qquad
    W^-.
\]
\end{definitionbox}

\begin{conceptbox}
Weil $W^\pm$ elektrische Ladung tragen, koennen sie Ladung zwischen Fermionen
uebertragen.

Beispiel:
\[
    d \rightarrow u+W^-.
\]

Das Fermion wird um eine Ladungseinheit positiver.
Das $W^-$ traegt die negative Ladung weg.
\end{conceptbox}

\begin{examtipbox}
$W$-Boson:
\[
    \text{geladen}
    \quad\Rightarrow\quad
    \text{Ladungsaenderung am Fermion moeglich}.
\]
\end{examtipbox}

\section{Higgs-Mechanismus}

In der elektroschwachen Theorie erhalten $W$- und $Z$-Bosonen ihre Masse durch den
Higgs-Mechanismus.

\begin{definitionbox}
Der \emph{Higgs-Mechanismus} ist der Mechanismus, durch den die $W$- und $Z$-Bosonen
Masse erhalten, waehrend das Photon masselos bleibt.
\end{definitionbox}

\begin{conceptbox}
Ohne Masse waere die schwache Wechselwirkung langreichweitig.
Die beobachtete kurze Reichweite der schwachen Wechselwirkung zeigt, dass $W$ und
$Z$ massiv sind.

Das Photon bleibt masselos.
Deshalb bleibt die elektromagnetische Wechselwirkung langreichweitig.
\end{conceptbox}

\begin{notebox}
Das Higgs-Boson ist die beobachtbare Anregung des Higgs-Feldes.
In diesem Kapitel ist vor allem wichtig:
Der Higgs-Mechanismus erklaert qualitativ, warum $W$ und $Z$ massiv sind und das
Photon nicht.
\end{notebox}

\section{Schwache Wechselwirkung und Erhaltungssaetze}

\begin{center}
\begin{tabular}{|p{5.0cm}|p{7.5cm}|}
\hline
Groesse & Verhalten in schwachen Prozessen \\
\hline
Energie und Impuls & erhalten \\
\hline
elektrische Ladung & erhalten \\
\hline
Baryonenzahl & in Standardprozessen erhalten \\
\hline
Leptonenzahl & in einfachen Standardprozessen erhalten \\
\hline
Quark-Flavor & kann sich aendern \\
\hline
Strangeness & kann verletzt werden \\
\hline
Paritaet & verletzt \\
\hline
\end{tabular}
\end{center}

\begin{examtipbox}
Die schwache Wechselwirkung ist speziell, weil sie Flavor aendern und Paritaet
verletzen kann.

Aber elektrische Ladung bleibt trotzdem immer erhalten.
\end{examtipbox}

\section{Indikatoren fuer schwache Prozesse}

\begin{conceptbox}
Ein Prozess ist wahrscheinlich schwach, wenn:
\begin{itemize}
    \item ein Neutrino beteiligt ist,
    \item ein $W^\pm$- oder $Z^0$-Boson beteiligt ist,
    \item ein Quark-Flavor wechselt,
    \item Strangeness nicht erhalten ist,
    \item Paritaet verletzt wird,
    \item die Lebensdauer relativ lang ist im Vergleich zu starken Zerfaellen.
\end{itemize}
\end{conceptbox}

\begin{examtipbox}
Typischer Pruefungsweg:
\begin{enumerate}
    \item Reaktion aufschreiben.
    \item Ladung pruefen.
    \item Baryonenzahl pruefen.
    \item Leptonenzahl pruefen.
    \item Flavor und Strangeness pruefen.
    \item Wenn Flavor verletzt wird: schwache Wechselwirkung.
\end{enumerate}
\end{examtipbox}

\section{Vergleich der Wechselwirkungen}

\begin{center}
\begin{tabular}{|p{3.6cm}|p{3.4cm}|p{3.5cm}|p{3.6cm}|}
\hline
Wechselwirkung & Austauschteilchen & Reichweite & typische Eigenschaft \\
\hline
elektromagnetisch & $\gamma$ & unendlich & wirkt auf elektrische Ladung \\
\hline
stark & $g$ & Hadronenskala & wirkt auf Farbladung \\
\hline
schwach & $W^\pm$, $Z^0$ & sehr kurz & Flavorwechsel, Neutrinos \\
\hline
\end{tabular}
\end{center}

\begin{conceptbox}
Elektromagnetische und starke Wechselwirkung erhalten Quark-Flavor.

Die schwache Wechselwirkung kann Quark-Flavor aendern.
Das macht sie fuer Teilchenzerfaelle besonders wichtig.
\end{conceptbox}

\section{Beispiel: schwacher Strange-Zerfall}

Seltsame Teilchen zeigen besonders gut den Unterschied zwischen starker Erzeugung
und schwachem Zerfall.

\begin{examplebox}
Betrachte:
\[
    \Lambda^0 \rightarrow p+\pi^-.
\]

Quarkinhalte:
\[
    \Lambda^0=uds,
\]
\[
    p=uud,
\]
\[
    \pi^-=d\overline{u}.
\]

Strangeness links:
\[
    S=-1.
\]

Strangeness rechts:
\[
    S=0.
\]

Also:
\[
    \Delta S=+1.
\]
Der Prozess ist schwach.
\end{examplebox}

\begin{conceptbox}
Auf Quarkebene kann man den Flavorwechsel schematisch als
\[
    s \rightarrow u+W^-
\]
auffassen.
Das $W^-$ erzeugt dann passende Endteilchen.
\end{conceptbox}

\section{Beispiel: Pionzerfall}

Geladene Pionen zerfallen schwach.

\begin{formulabox}
\[
    \pi^- \rightarrow \mu^- + \overline{\nu}_\mu.
\]
\end{formulabox}

\begin{conceptbox}
Quarkinhalt:
\[
    \pi^- = d\overline{u}.
\]

Das Quark-Antiquark-Paar kann ueber ein virtuelles $W^-$ annihilieren:
\[
    d+\overline{u}\rightarrow W^-,
\]
\[
    W^- \rightarrow \mu^-+\overline{\nu}_\mu.
\]
\end{conceptbox}

\begin{examtipbox}
Ein Hadron kann schwach in Leptonen zerfallen, wenn ein virtuelles $W$-Boson
vermittelt.
\end{examtipbox}

\section{Mini-Zusammenfassung}

\begin{notebox}
Die wichtigsten Punkte dieses Kapitels sind:
\begin{itemize}
    \item Die schwache Wechselwirkung vermittelt Beta-Zerfaelle, Neutrinoprozesse
    und Flavorwechsel.
    \item Die Austauschteilchen sind:
    \[
        W^+,\quad W^-,\quad Z^0.
    \]
    \item $W^\pm$ vermitteln den geladenen schwachen Strom.
    \item $Z^0$ vermittelt den neutralen schwachen Strom.
    \item $W^\pm$ koennen die Ladung und den Flavor eines Fermions aendern.
    \item $Z^0$ ist neutral und aendert die elektrische Ladung nicht.
    \item Beta-Minus-Zerfall:
    \[
        n \rightarrow p+e^-+\overline{\nu}_e.
    \]
    \item Auf Quarkebene:
    \[
        d \rightarrow u+W^-.
    \]
    \item Beta-Plus-Zerfall:
    \[
        p \rightarrow n+e^++\nu_e.
    \]
    \item Auf Quarkebene:
    \[
        u \rightarrow d+W^+.
    \]
    \item Neutrinos wechselwirken schwach.
    \item Die schwache Wechselwirkung verletzt Paritaet.
    \item Die Fermi-Theorie ist die niederenergetische Naeherung der schwachen
    Wechselwirkung.
    \item Elektromagnetische und schwache Wechselwirkung sind in der elektroschwachen
    Theorie vereinheitlicht.
    \item Der Higgs-Mechanismus erklaert qualitativ die Masse von $W$ und $Z$ und
    das masselose Photon.
\end{itemize}
\end{notebox}

\section{Pruefungscheck}

\begin{examtipbox}
Nach diesem Kapitel solltest du folgende Fragen beantworten koennen:
\begin{enumerate}
    \item Was ist die schwache Wechselwirkung?
    \item Welche Austauschteilchen vermitteln die schwache Wechselwirkung?
    \item Was ist der geladene schwache Strom?
    \item Was ist der neutrale schwache Strom?
    \item Warum ist die schwache Wechselwirkung kurzreichweitig?
    \item Welche Rolle spielen $W^+$ und $W^-$?
    \item Welche Rolle spielt das $Z^0$?
    \item Wie sieht der Beta-Minus-Zerfall auf Nukleonenebene aus?
    \item Wie sieht der Beta-Minus-Zerfall auf Quarkebene aus?
    \item Wie sieht der Beta-Plus-Zerfall auf Nukleonenebene aus?
    \item Wie sieht der Beta-Plus-Zerfall auf Quarkebene aus?
    \item Was passiert beim Elektroneneinfang?
    \item Warum sind Neutrinos schwer nachzuweisen?
    \item Was ist die inverse Beta-Reaktion?
    \item Warum kann die schwache Wechselwirkung Flavor aendern?
    \item Was beschreibt die CKM-Matrix?
    \item Was bedeutet Paritaetsverletzung?
    \item Was ist die Fermi-Theorie?
    \item Warum ist die Fermi-Theorie eine Naeherung?
    \item Was bedeutet elektroschwache Vereinheitlichung?
    \item Warum ist das Photon masselos, aber $W$ und $Z$ massiv?
    \item Woran erkennt man in Reaktionen schwache Prozesse?
\end{enumerate}
\end{examtipbox}